% !TeX root = ../../Book3.tex
% 纲要标题：### 23.3 光滑与平展
% 纲要位置：计划纲要.md，第 2717 行。

\section{光滑与平展}
\label{b3:sec:2303}

可逆 Jacobian 子式允许独立修正一部分坐标，其余坐标仍可自由变化。标准模型使这一计算成为局部方程组，而平坦性与唯一提升还需要各自的代数论证。

% 纲要内容（仅作写作计划，尚非教材正文）：
%
% * 标准光滑模型的自由变量、真实关系与 Jacobian 子式；Nakayama 排除未列出的关系
% * 采用 Noether 局部环的极大理想进形式光滑判据，核对离散提升到连续提升；以有限系数子环推广平坦性
% * 对角幂等元的模分裂及平坦性转移；有限展示下平展的等价刻画，平坦性消去额外商理想
% * 平方覆盖的函数域次数阻止 Zariski 开同构；固定系数坐标取幂的自由秩与代数闭域纤维长度，比较绝对 Frobenius 并区别点集和无穷小性质
%
% ---
%

\subsection{光滑映射的标准局部模型}
\label{b3:subsec:2303:01}

\begin{definition}\label{b3:def:standard-smooth-model}
设 \(A\) 为交换含幺环，\(n\ge0\)、\(0\le r\le n\)，并取 \(f_1,\ldots,f_r,g\in A[X_1,\ldots,X_n]\)。令
\[
B=\bigl(A[X_1,\ldots,X_n]/(f_1,\ldots,f_r)\bigr)_g.
\]
如果 Jacobian 矩阵 \((\partial f_i/\partial X_j)\) 的某个 \(r\) 阶子式 \(\Delta\) 在 \(B\) 中为单位，就称这个展示为一个\term[biaozhun-guanghua-moxing]{标准光滑模型}[standard smooth presentation]。当 \(r=n\) 时称为\term[biaozhun-pingzhan-moxing]{标准平展模型}[standard étale presentation]。\(r=0\) 时空子式取为 \(1\)，得到多项式环的局部化。
\end{definition}

可逆的 \(r\) 阶子式允许用 \(r\) 个坐标修正 \(r\) 条独立关系，余下 \(n-r\) 个坐标仍能自由变化。下一个命题将在微分模中把这个数表现为基向量的数目；\(r=n\) 时，全部变量都参与修正，不再留下无穷小变化方向。

\begin{proposition}\label{b3:prop:standard-smooth-lifting}
设 \(A\) 为交换含幺环，\(n\ge0\)、\(0\le r\le n\)，并取 \(f_1,\ldots,f_r,g\in A[X_1,\ldots,X_n]\)。令 \(B=(A[X_1,\ldots,X_n]/(f_1,\ldots,f_r))_g\)。若 Jacobian 矩阵 \((\partial f_i/\partial X_j)\) 的一个 \(r\) 阶子式在 \(B\) 中可逆，则 \(A\to B\) 光滑。若可逆子式由前 \(r\) 列组成，则 \(\Omega_{B/A}\) 自由，以 \(\mathrm{d}X_{r+1},\ldots,\mathrm{d}X_n\) 为基。特别地，\(r=n\) 时 \(A\to B\) 形式平展。
\end{proposition}
\begin{proof}
重排变量后，可设给定可逆子式由前 \(r\) 列组成。有限展示显然成立，因为逆元可增加一个生成元及关系 \(gY-1\)。对 \(u:B\to C/J\)、\(J^2=0\)，任取变量像的提升 \(c_i\)。\(g(c)\) 与 \(\Delta(c)\) 都是单位，因为它们模 \(J\) 是单位。将后 \(n-r\) 个变量保持不动，令前 \(r\) 个变量的修正 \(h_j\in J\) 解线性方程
\[
\sum_{j=1}^r\frac{\partial f_i}{\partial X_j}(c)h_j=-f_i(c),
\qquad 1\le i\le r.
\]
系数矩阵可逆，故有唯一这样的修正。平方零 Taylor 公式使全部方程精确成立，局部化所需的单位也仍是单位，故得到提升。

微分模的关系矩阵恰为这个 Jacobian 矩阵。当 \(0<r<n\) 时，将它写为 \((J_0\mid J_1)\)，其中 \(J_0\) 是可逆的前 \(r\) 列。将关系左乘 \(J_0^{-1}\)，得到
\[
\begin{pmatrix}\mathrm{d}X_1\\\vdots\\\mathrm{d}X_r\end{pmatrix}
=-J_0^{-1}J_1
\begin{pmatrix}\mathrm{d}X_{r+1}\\\vdots\\\mathrm{d}X_n\end{pmatrix}.
\]
于是前 \(r\) 个微分均由其余微分表示。若这些关系的线性组合只涉及其余微分，则其前 \(r\) 个系数都为零；归一后的左侧系数矩阵是单位阵，故这个线性组合的系数只能全为零。因此其余微分之间没有非零关系，证明所述自由性。\(r=0\) 时没有关系，全部 \(\mathrm{d}X_j\) 成为自由基；\(r=n>0\) 时可逆系数矩阵使全部微分为零。因此 \(r=n\) 时总有 \(\Omega_{B/A}=0\)，由形式非分歧判据及已证的形式光滑性，得到形式平展。
\end{proof}

标准模型不仅给出一类例子。有限展示使任意光滑代数都能在每一点附近写成这样的形式。取有限多项式展示 \(B=P/I\)，证明须从关系模的局部基选出真正的关系 \(f_i\)，再用分裂的微分映射得到可逆 Jacobian 子式，最后用 Nakayama 引理证明这些 \(f_i\) 在相应局部化中确实生成整个 \(I\)。Jacobian 检验关系的一阶独立性，最后一步才排除尚未列出的实际关系。

\begin{theorem}[光滑映射的局部标准形]\label{b3:thm:smooth-local-standard-form}
设 \(A\) 为交换含幺环，\(B\) 为有限展示交换 \(A\) 代数。则 \(A\to B\) 光滑，当且仅当对每个 \(\mathfrak q\in\operatorname{Spec}B\)，存在 \(b\notin\mathfrak q\)，使 \(B_b\) 在 \(A\) 上具有标准光滑模型。
\end{theorem}
\begin{proof}
先设 \(B\) 形式光滑。写 \(B=P/I\)，其中 \(P=A[X_1,\ldots,X_n]\)，\(I\) 有限生成。由\cref{b3:thm:formal-smooth-conormal-splitting}，\(I/I^2\) 是 \(B^n\) 的直和项，故为有限生成投射模。在给定 \(\mathfrak q\) 的某个主开邻域上，它自由；从 \(I\) 的一组有限生成元中可选取 \(f_1,\ldots,f_r\)，使其像成为这一邻域上的基。更具体地，先在剩余域上选出基，再用局部自由性及相应行列式非零，缩小邻域使它们确为基。

设 \(Q\subseteq P\) 为 \(\mathfrak q\) 的逆像。因为 \(I/I^2\to B^n\) 分裂单射，\(\mathrm{d}f_1,\ldots,\mathrm{d}f_r\) 在剩余域上仍独立，所以它们的 Jacobian 矩阵有一个 \(r\) 阶子式 \(\Delta\notin Q\)。还要证明这些 \(f_i\) 在邻域上生成 \(I\)，而不只是模 \(I^2\) 生成。令 \(F=(f_1,\ldots,f_r)\)。在局部环 \(P_Q\) 中，
\[
I_Q=F_Q+I_Q^2.
\]
有限生成模 \(N=I_Q/F_Q\) 因而满足 \(N=I_QN\)。由于 \(I_Q\subseteq QP_Q\)，Nakayama 引理给出 \(N=0\)。模 \(I/F\) 有限生成，所以它在 \(Q\) 处为零意味着存在一个 \(h\notin Q\)，使 \(I_h=F_h\)：对有限组生成元分别选择零化它们的分母，再取乘积即可。因此
\[
B_{\bar h\bar\Delta}
\cong\bigl(P/(f_1,\ldots,f_r)\bigr)_{h\Delta},
\]
这就是包含 \(\mathfrak q\) 的标准光滑邻域。这里没有假设 \(A\) 为 Noether 环；所用有限性来自 \(I\) 有限生成。

反之，设这样的邻域覆盖 \(\operatorname{Spec}B\)。仍取固定的有限展示 \(B=P/I\)。微分模 \(\Omega_{B/A}\) 有限展示，且由\cref{b3:prop:standard-smooth-lifting}在这些邻域上有限自由，故由\cref{b3:thm:finite-presentation-flat-projective}为有限生成投射模。另一方面，在每个邻域中 \(B_b\) 形式光滑，关系模判据对局部化的多项式环仍适用，故固定映射 \(I/I^2\to B^n\) 在该邻域上单射。其核处处局部化为零，所以全局单射。于是
\(0\to I/I^2\to B^n\to\Omega_{B/A}\to0\)
正合。因末项投射，此列分裂，第一箭头有左逆。再用关系模判据，得到 \(B\) 形式光滑。
\end{proof}

\subsection{平展映射的平坦性与非分歧性}
\label{b3:subsec:2303:02}

第14章对首一单变量模型，已经用商环的自由基直接证明平坦性。一般标准光滑模型不一定以这种单变量形式出现；从提升性质推出平坦性，需要一个更深的局部结果。对 Noether 局部环同态 \((A,\mathfrak m)\to(B,\mathfrak n)\)，松村英之\cite[\S28, pp. 213--215; Theorem 28.9, p. 222]{Matsumura1989CommutativeRingTheory}将 \(\mathfrak n\)-进形式光滑性刻画为平坦性与闭纤维的相应形式光滑性。该等价定理的证明较深，松村书也转引其原始证明。这里采用其中推出平坦性的方向；下面的概要核对本书的形式光滑定义如何满足其拓扑假设，所采用的局部平坦性判据本身不在此推导。

\begin{theorem}[Noether 局部形式光滑蕴含平坦]\label{b3:thm:local-formal-smooth-flat-external}
设 \((A,\mathfrak m)\to(B,\mathfrak n)\) 为交换 Noether 局部环之间的局部同态。若 \(B\) 在 \(A\) 上形式光滑，则 \(B\) 是平坦 \(A\) 模。
\end{theorem}

\begin{proofsketch}
先说明离散提升性质为什么满足松村定理的拓扑假设。令 \(C\) 为交换 \(A\) 代数，\(J\subseteq C\) 为平方零理想，并给定到离散环的连续映射 \(u:B\to C/J\)，即 \(u(\mathfrak n^v)=0\) 对某个 \(v\ge1\) 成立。本书的形式光滑性给出提升 \(w:B\to C\)，而
\[
w(\mathfrak n^v)\subseteq J,
\qquad w(\mathfrak n^{2v})\subseteq J^2=0.
\]
因此提升也连续。这正是 \(B\) 在 \(A\) 上 \(\mathfrak n\)-进形式光滑的含义。

现在记 \(k=A/\mathfrak m\)、\(B_0=B/\mathfrak mB\)、\(\mathfrak n_0=\mathfrak n/\mathfrak mB\)。松村的局部判据断言，\(B\) 在 \(A\) 上 \(\mathfrak n\)-进形式光滑，当且仅当它在 \(A\) 上平坦，且闭纤维 \(B_0\) 在 \(k\) 上 \(\mathfrak n_0\)-进形式光滑\cite[Theorem 28.9, p. 222]{Matsumura1989CommutativeRingTheory}。前一段已验证其左侧条件，故可用这一判据得到平坦性。闭纤维条件是该判据给出的另一个结论，无须另加到本定理的假设中。
\end{proofsketch}

\begin{theorem}\label{b3:thm:smooth-flat}
设 \(A\to B\) 为交换含幺环之间的光滑同态，则 \(B\) 在 \(A\) 上平坦。
\end{theorem}
\begin{proof}
标准模型只用到有限多个系数。先把这些系数放到一个 Noether 子环上，逐点应用前面的局部平坦性定理，再将得到的平坦模型基变换回 \(A\)，就能处理任意底环。

先处理标准模型 \(B=(A[X_1,\ldots,X_n]/(f_1,\ldots,f_r))_g\)，设其可逆子式为 \(\Delta\)。再使 \(\Delta\) 可逆不改变 \(B\)，故可把局部化元素写成 \(g\Delta\)。令 \(A_0\) 为由 \(\mathbb Z\) 在 \(A\) 中的像以及 \(f_i,g\) 的有限多个系数生成的子环。它是有限型 \(\mathbb Z\) 代数，因而为 Noether 环。\(\Delta\) 由 \(f_i\) 求导、取行列式得到，其系数也在 \(A_0\) 中。定义
\[
B_0=\bigl(A_0[X_1,\ldots,X_n]/(f_1,\ldots,f_r)\bigr)_{g\Delta}.
\]
这是 Noether 环；\(\Delta\) 已在此模型中显式可逆，且有限展示与局部化的基变换给出 \(B\cong A\otimes_{A_0}B_0\)。标准模型的提升证明给出 \(A_0\to B_0\) 形式光滑。对每个素理想 \(\mathfrak q\subseteq B_0\)，令 \(\mathfrak p=\mathfrak q\cap A_0\)。先沿 \(A_0\to(A_0)_{\mathfrak p}\) 作基变换，得到形式光滑的同态
\[
(A_0)_{\mathfrak p}\longrightarrow
C=(A_0)_{\mathfrak p}\otimes_{A_0}B_0.
\]
因为 \(A_0\setminus\mathfrak p\) 与 \(\mathfrak q\) 不相交，\(\mathfrak q\) 在 \(C\) 中扩张为素理想 \(\mathfrak qC\)。再对目标环局部化，得到仍然形式光滑的局部同态
\[
(A_0)_{\mathfrak p}\longrightarrow
C_{\mathfrak qC}\cong(B_0)_{\mathfrak q}.
\]
前一局部定理因此给出 \((B_0)_{\mathfrak q}\) 在 \((A_0)_{\mathfrak p}\) 上平坦；而 \((A_0)_{\mathfrak p}\) 是平坦 \(A_0\) 模，故 \((B_0)_{\mathfrak q}\) 在 \(A_0\) 上也平坦。

这足以证明 \(B_0\) 在 \(A_0\) 上平坦。事实上，对任意 \(A_0\) 模单射 \(U\to V\)，张量后映射的核是一个 \(B_0\) 模；把它在每个 \(\mathfrak q\) 处局部化，得到与平坦模 \((B_0)_{\mathfrak q}\) 张量后的核，均为零。由模的局部零检测，这个核为零。最后由平坦性在基变换下保持，\(B=A\otimes_{A_0}B_0\) 在 \(A\) 上平坦。

一般光滑代数由\cref{b3:thm:smooth-local-standard-form}在目标上被这样的标准模型覆盖。同样对张量映射的核逐个邻域作局部化，就得到全局平坦性。
\end{proof}

为了从“平坦且微分为零”得到唯一提升，还须把底环上的平坦性转到一个候选标准平展环上。对有限型 \(A\) 代数 \(C\)，微分模为零使对角理想 \(K=\ker(C\otimes_AC\to C)\) 满足 \(K/K^2=0\)。有限型保证 \(K\) 有限生成，于是 \(K=K^2\) 使它由一个幂等元生成。其互补幂等元将任意 \(C\) 模 \(M\) 从 \(C\otimes_AM\) 中分裂出来，正是下面引理转移平坦性的机制。

\ProseParagraphBreak
先在 \(C=A\times A\) 中看这个分裂。以 \(\mu:C\otimes_AC\to C\) 为乘法映射，令 \(e_1=(1,0)\)、\(e_2=(0,1)\)，并取
\[
z=e_1\otimes e_1+e_2\otimes e_2
\]
则 \(\mu(z)=1\) 使 \(m\mapsto e_1\otimes e_1m+e_2\otimes e_2m\) 成为乘法映射 \(C\otimes_AM\to M\) 的截面，而 \((c\otimes1)z=(1\otimes c)z\) 保证这个截面为 \(C\) 线性。两个等式分别完成复合为恒等与标量作用相容这两件事；下面构造一般情形中的同一个 \(z\)。

\begin{lemma}\label{b3:lem:unramified-diagonal-splitting}
设 \(A\) 为交换含幺环，\(C\) 为有限型交换 \(A\) 代数，且 \(\Omega_{C/A}=0\)。若 \(M\) 为一个 \(C\) 模并且在 \(A\) 上平坦，则 \(M\) 在 \(C\) 上也平坦。
\end{lemma}
\begin{proof}
令 \(D=C\otimes_AC\)，\(\mu:D\to C\) 为乘法，\(K=\ker\mu\)。若 \(c_1,\ldots,c_l\) 生成 \(C\)，则 \(K\) 由有限个 \(c_i\otimes1-1\otimes c_i\) 生成：模去这些差以后，两份 \(C\) 的像相同，商正是 \(C\)。映射 \(c\mapsto c\otimes1-1\otimes c\bmod K^2\) 为导子，并由其泛性质给出
\(K/K^2\cong\Omega_{C/A}\)。也可直接用这些生成元及乘积的求导关系验证这个同构。因此 \(K=K^2\)。

有限生成且幂等的理想由一个幂等元生成。为说明这一点，取生成元列 \(a\)，由 \(K=K^2\) 写 \(a=Ea\)，其中矩阵 \(E\) 的各项属于 \(K\)。伴随矩阵给出 \(\det(1-E)K=0\)。写 \(\det(1-E)=1-e\)、\(e\in K\)，则 \((1-e)e=0\)，即 \(e^2=e\)，而每个 \(x\in K\) 满足 \(x=ex\)，所以 \(K=(e)\)。令 \(z=1-e\)，则
\[
\mu(z)=1,\qquad(c\otimes1-1\otimes c)z=0\quad(c\in C).
\]
写 \(z=\sum_i u_i\otimes v_i\)。映射
\[
s:M\longrightarrow C\otimes_AM,
\qquad m\longmapsto\sum_i u_i\otimes v_i m
\]
由第二个等式为 \(C\) 线性，且乘法映射 \(C\otimes_AM\to M\) 与它的复合为恒等映射。故 \(M\) 是 \(C\otimes_AM\) 的直和项。后者由基变换在 \(C\) 上平坦，直和项也平坦，结论成立。
\end{proof}

\begin{theorem}[平展的等价刻画]\label{b3:thm:etale-equivalent-characterizations}\label{b3:thm:closed-point-unramified-not-etale}
设 \(A\) 为交换含幺环，\(B\) 为有限展示交换 \(A\) 代数。下列条件等价：
\begin{equivalences}
\item\label{b3:item:etale-flat-unramified} \(B\) 在 \(A\) 上平坦，且 \(\Omega_{B/A}=0\)，即 \(A\to B\) 平展；
\item\label{b3:item:etale-formal} \(A\to B\) 形式平展；
\item\label{b3:item:etale-standard} 每个 \(\mathfrak q\in\operatorname{Spec}B\) 都有主开邻域，在该邻域上 \(B\) 具有标准平展模型。
\end{equivalences}
其中平坦性不能仅由微分消失代替：对任意域 \(k\)，\(A=k[t]\to B=A/(t)=k\) 满足 \(\Omega_{B/A}=0\)，是有限展示的非分歧映射，但不平坦、不形式光滑，因而不平展。
\end{theorem}
\begin{proof}
先从条件\itemref{b3:item:etale-formal}出发。形式平展包含形式光滑性，由局部标准形定理得到标准光滑邻域 \(D(b)\)，其中包含给定的 \(\mathfrak q\)。形式非分歧又使 \(\Omega_{B/A}=0\)，而标准模型给出
\[
0=\Omega_{B_b/A}\cong B_b^{\,n-r}.
\]
由于 \(\mathfrak q\in D(b)\)，环 \(B_b\) 有素理想，因而 \(B_b\ne0\)。非零环上的自由模只有在秩为零时才为零，故 \(n=r\)，模型标准平展，得到条件\itemref{b3:item:etale-standard}。

再设条件\itemref{b3:item:etale-standard}成立。标准平展邻域由局部标准形定理使 \(B\) 光滑，且微分模处处为零，因而全局为零。这给出形式光滑与形式非分歧，合起来即为条件\itemref{b3:item:etale-formal}。光滑蕴含平坦又给出条件\itemref{b3:item:etale-flat-unramified}。

最后从条件\itemref{b3:item:etale-flat-unramified}构造标准平展邻域。先用微分消失选出一组 Jacobian 可逆的关系，得到一个候选标准平展环；原代数仍可能是它的商，平坦性将消去这些尚未列出的关系。写 \(B=P/I\)、\(P=A[X_1,\ldots,X_n]\)，其中 \(I\) 有限生成。给定 \(\mathfrak q\)，关系正合列和 \(\Omega_{B/A}=0\) 说明 \(I/I^2\to B^n\) 满射。模去 \(\mathfrak q\) 后，这些生成元的梯度向量张成 \(\kappa(\mathfrak q)^n\)，从中选一组基，便得 \(n\) 个关系 \(f_1,\ldots,f_n\)，使其 Jacobian 行列式满足 \(\bar\Delta\notin\mathfrak q\)。令
\[
C=\bigl(P/(f_1,\ldots,f_n)\bigr)_\Delta,
\qquad B'=B_{\bar\Delta}=C/J.
\]
这里 \(J\) 由其余有限个关系生成，\(C\) 为标准平展模型，故 \(\Omega_{C/A}=0\)。\(B'\) 在 \(A\) 上平坦，因为它是平坦 \(A\) 代数 \(B\) 的局部化。由\cref{b3:lem:unramified-diagonal-splitting}，\(B'\) 在 \(C\) 上平坦。

用平坦 \(C\) 模 \(B'\) 张量单射 \(J\to C\)，得到单射 \(J/J^2\to B'\)；这个映射把 \(j\) 送到其在 \(C/J\) 中的零像，故本身为零，只能有 \(J/J^2=0\)，即 \(J=J^2\)。将所选 \(\mathfrak q\) 扩张到 \(B'\)，令 \(\mathfrak r\subseteq C\) 为其逆像，则 \(J\subseteq\mathfrak r\)。在 \(C_{\mathfrak r}\) 中，有限生成模 \(J_{\mathfrak r}\) 满足
\(J_{\mathfrak r}=J_{\mathfrak r}^2\)，由 Nakayama 引理为零。因此存在 \(h\notin\mathfrak r\)，使 \(J_h=0\)。于是 \(B'_h=C_h\)，这是包含给定点的标准平展邻域。这样条件\itemref{b3:item:etale-flat-unramified}也推出条件\itemref{b3:item:etale-standard}。

最后的闭点反例已在\cref{b3:exm:unramified-not-etale}证明：相对导子全为零，乘 \(t\) 的单射张量 \(k\) 后成为零映射，而 \(k\to k[t]/(t^2)\) 的 \(k[t]\) 代数提升不能使 \(t\) 的两个像相容。在上面的候选模型构造中，这个例子取 \(C=A=k[t]\)、\(J=(t)\)，仍有 \(J/J^2\cong k\ne0\)。因此微分消失本身没有删掉这条关系，平坦性正是所缺的条件。
\end{proof}

这样便得到了第14章提出的定义比较。有限展示在选择有限组关系、把一点处的消失扩大到邻域，以及形成标准模型时都有实际作用；平坦性则在最后一个证明中消去了多余的商理想。只要求 \(\Omega_{B/A}=0\) 不足以做到这一点，\(k[t]\to k\) 正是已经算过的反例。

\ProseParagraphBreak
对交换含幺环同态 \(A\to B\)，上述形式性质与有限性条件的关系汇总于\cref{b3:tab:formal-geometric-finiteness}。其中的局部模型均指在 \(\operatorname{Spec}B\) 的每一点附近成立的展示。

\begin{table}[htbp]
\centering
\caption{形式性质、有限性条件与局部模型}
\label{b3:tab:formal-geometric-finiteness}
\begin{tabularx}{\textwidth}{@{}l l X@{}}
\toprule
形式性质 & 附加有限性 & 所得性质与局部模型 \\
\midrule
形式光滑 & 有限展示 & 光滑，提升存在；局部具有标准光滑模型，且必平坦 \\
形式非分歧 & 有限型 & 非分歧，提升至多唯一；\(\Omega_{B/A}=0\)，但不一定平坦 \\
形式平展 & 有限展示 & 平展，提升存在且唯一；局部具有标准平展模型 \\
\bottomrule
\end{tabularx}
\end{table}

平展一行也可以用有限展示、平坦及 \(\Omega_{B/A}=0\) 表示。这里的有限展示与有限型都是指 \(B\) 作为 \(A\) 代数的有限性条件。

\subsection{局部同胚直觉与平展的代数条件}
\label{b3:subsec:2303:03}

在复数几何中，可逆 Jacobian 容易使人联想到逆函数定理。这个直觉有助于理解“相对零维而没有无穷小变化”，却不能把平展直接定义成 Zariski 拓扑中的局部同构，更不能只检查点集映射。

\begin{example}\label{b3:exm:etale-not-zariski-isomorphism}
设 \(A=\mathbb C[s,s^{-1}]\)、\(B=\mathbb C[t,t^{-1}]\)，映射由 \(s\mapsto t^2\) 给出。写成
\(B\cong A[T]/(T^2-s)\)，其中 \(T\) 自动可逆，导数 \(2T\) 为单位。故此映射平展，且 \(B\) 为秩二自由 \(A\) 模。

\ProseParagraphBreak
然而不能在 \(\operatorname{Spec}B\)、\(\operatorname{Spec}A\) 中分别取非空 Zariski 开集，使给定谱映射的限制成为同构。这样的相容同构必诱导函数域同构，迫使 \(\mathbb C(s)\to\mathbb C(t)\)、\(s\mapsto t^2\) 为同构。但在函数域 \(\mathbb C(s)\) 的 \(s=0\) 离散赋值 \(v_0\) 下，\(v_0(s)=1\)，而非零平方的赋值必为偶数，故 \(T^2-s\) 在 \(\mathbb C(s)\) 上不可约，扩张次数为二。给定映射在函数域上已经有次数二的障碍，缩小非空开集不能消去它。这里用的是函数域上的赋值，\(s=0\) 不是 \(\operatorname{Spec}A\) 中的点；复解析拓扑中可以局部选定平方根，采用的是另一种邻域。
\end{example}

\begin{example}\label{b3:exm:frobenius-homeomorphism-not-etale}
设 \(k\) 为特征 \(p>0\) 的代数闭域，考虑 \(k[s]\to k[t]\)、\(s\mapsto t^p\)。目标以 \(1,t,\ldots,t^{p-1}\) 为自由基，所以映射有限且平坦。在素谱上，它把泛点送到泛点，并把闭点 \((t-a)\) 送到 \((s-a^p)\)；代数闭域上 \(a\mapsto a^p\) 为双射。有限整映射是闭映射，因此这个素谱映射为同胚。

\ProseParagraphBreak
这里保持 \(k\) 中系数不变，只对坐标取 \(p\) 次幂。它与对每个元素 \(r\) 都取 \(r^p\) 的绝对 Frobenius 一般不同；后者同时改变系数，且在素谱上逐点恒等。附录 C 的\secref{b3:sec:C02}将明确比较这两种环同态。

\ProseParagraphBreak
其相对微分模却为
\[
\Omega_{k[t]/k[s]}
\cong k[t]\,\mathrm{d}t/(\mathrm{d}(t^p))=k[t]\,\mathrm{d}t\ne0.
\]
故它不是非分歧映射，也不是平展映射。

\ProseParagraphBreak
进一步，在闭点 \(s=a\) 上，由代数闭性可唯一选取 \(b\in k\) 使 \(b^p=a\)。纤维代数为
\[
k[t]\otimes_{k[s]}k
\cong k[t]/(t^p-a)
=k[t]/((t-b)^p)
\cong k[\varepsilon]/(\varepsilon^p).
\]
它以 \(1,\varepsilon,\ldots,\varepsilon^{p-1}\) 为 \(k\) 基，维数为 \(p\)。这个局部 Artin 代数的剩余域仍是 \(k\)，所以其组成列长度也为 \(p\)，等于原有限自由扩张的秩。各闭纤维都只有一个底层点，却带有非零幂零元 \(\varepsilon\)；这种厚度与平坦性共存，而这份纤维代数在 \(k\) 上的微分模仍为非零的 \(k[\varepsilon]/(\varepsilon^p)\,\mathrm{d}\varepsilon\)。点集及普通拓扑看不到这部分无穷小信息。
\end{example}

因此，对给定代数映射，平展性应当由有限展示、平坦性、微分模，或由已经证明等价的标准模型来检验。只知道两个环都正则，或者两个素谱同胚，都不够。
